% ---------------------------------------------------------------------
\clearpage
\section{Introducción}\label{cap:introduccion}

\subsection{Motivación}

El síndrome de Brugada es una enfermedad eléctrica del corazón de origen genético, la cual fue descrita por los hermanos Brugada en el año 1992, la cual predispone a la fibrilación ventricular y a la muerte súbita en personas jóvenes sin que estas cuenten con alguna cardiopatía estructural. En muchas ocasiones, la primera manifestación es el propio evento arrítmico. Por eso, es muy importante identificar el patrón electrocardiográfico en los ECG rutinarios, para que de esta forma los médicos sean capaces de estudiar al paciente y a su familia, evitar los fármacos o situaciones que puedan llegar a desencadenar arritmias, e incluso llegar a indicar un desfibrilador implantable.

El diagnóstico para esta cardiopatía se basa en un patrón característico en las derivaciones precordiales derechas V1 y V2, caracterizado por una elevación del segmento ST con un descenso convexo, seguido de una onda T negativa. Sin embargo, este patrón no siempre es el mismo: cuando aparece puede llegar a ser bastante sutil y puede llegar a confundirse con algunas variantes benignas como puede llegar a ser el bloqueo incompleto de la rama derecha. En el conjunto de datos que usaré dentro de este trabajo, de los 76 pacientes de la clase positiva, solo 23 de ellos tienen marcado el ECG basal como patológico. Es debido a esta gran dificultad para su diagnóstico que el uso del aprendizaje automático puede llegar a ser muy interesante.

La mayoría de los trabajos publicados acerca de esta enfermedad se centran en el modelo, pero en este problema, como en la mayoría de los problemas reales que nos encontramos dentro del campo de la Inteligencia de Negocio, el hecho de obtener o no un buen resultado depende de la calidad de los datos que vayamos a utilizar. Por eso, partimos de un ECG de 12 derivaciones de \SI{12}{\second} a \SI{100}{\hertz}, el cual contiene \num{14400} valores, que tienen ruido, pueden tener deriva en la línea de base, latidos anómalos, y el patrón que lo define dura apenas \SI{200}{\milli\second}, para tratar de convertirlo en unas pocas variables que sean fiables y que tengan de verdad un significado clínico. Es esta parte del KDD la que más esfuerzo exige y la que de verdad condiciona que cualquier modelo que entrenemos más tarde sea capaz de aprender.

\subsection{Objetivos}\label{sec:objetivos}

El objetivo general es el de construir y validar un proceso completo y reproducible capaz de transformar las señales ECG crudas del \emph{dataset} Brugada-HUCA en un conjunto de datos tabular listo para poder ser analizado. Los objetivos se corresponden con las siguientes fases del proceso:

\begin{enumerate}[label=\textbf{O\arabic*}, leftmargin=2.4em]
  \item \textbf{Proceso ETL y control de calidad}: extraer los registros, integrar estos con sus metadatos y detectar registros que estén duplicados, que tengan saturación, ruido o deriva.
  \item \textbf{Preprocesado de la señal}: filtrado, detección de latidos, segmentación y alineamiento, comparando varias alternativas para cada paso.
  \item \textbf{Selección de instancias}: rechazar los latidos que sean anómalos y construir un latido medio robusto por paciente.
  \item \textbf{Extracción y selección de características}: definir variables clínicas sobre el latido medio y comparar métodos de selección de características por rendimiento y coste.
  \item \textbf{Visualización}: explorar la propia estructura de los datos, usando técnicas de reducción de dimensionalidad, y cuantificar lo que están mostrando.
  \item \textbf{Modelo de referencia}: medir el valor que de verdad proporcionan estos datos, entrenando un clasificador sencillo, usando un protocolo de evaluación en el que no se produzca ninguna fuga de datos.
\end{enumerate}

\subsection{Contribuciones}

Además de cubrir los objetivos que acabo de nombrar, el trabajo también aporta:

\begin{itemize}
  \item Un proceso reproducible, que se encuentra en un repositorio abierto, con un comando por análisis, semillas fijas y una huella criptográfica del conjunto de datos.
  \item Un protocolo de evaluación sin fuga de información: todo lo que se ajusta con datos (normalización, selección de variables, hiperparámetros y umbral de decisión) se calcula dentro de cada partición de entrenamiento.
  \item Una comparación de métodos de selección de características de dos familias distintas (filtros y embebidos), evaluados sin sesgo de selección.
  \item Una confirmación cuantitativa de que la forma del segmento ST (su pendiente) discrimina mucho mejor que su elevación absoluta. No es una idea nueva, ya que coincide con criterios clínicos como el índice de Corrado, pero aquí se confirma con datos públicos y con métodos que no partían de ese conocimiento.
\end{itemize}
